\documentclass[oneside,a4paper,11pt,explicit]{book}
\usepackage[utf8]{inputenc}
\usepackage[sedes]{jamali}

%% INSERT YOUR PACKAGES HERE
\usepackage{lipsum}

%% TITLE PAGE OPTIONS (only works with \maketitle is called in the document)
\title{Judul Buku}
\subtitle{Template Buku di Penerbit Jala Marwita Lestari}
\date{\today}
\author{Bambang Purnomosidi D. P.}
%\professor{We Don't Know (yet)}

%% DOCUMENT
\begin{document}
\pagenumbering{roman}
\maketitle
\tableofcontents
\thispagestyle{empty}
\listoffigures
\listoftables
\lstlistoflistings
\newpage
\pagenumbering{arabic}

\chapter{Contoh Bab}
\jamalibox{{\bf Catatan:} Ini tulisan catatan awal bab. Silakan diisikan dengan berbagai hal yang dianggap penting untuk isi bab yang bersangkutan}
\lipsum[1]

\section{Ini contoh sub-bab}
\lipsum[2]

Contoh formula dalam teks: $f(x) = x^2$. Setelah itu, berikut ini adalah contoh formula dalam bagian tersendiri:

\begin{equation}
\ddot{x}+k\dot{x}+x^3=B\cos t,
\end{equation}
dengan $k=0.1$ dan $B=11$. 

Berikut adalah contoh pertanyaan:

\begin{questions}
\question Apakah anda bisa melihat pertanyaan 1 ini?
\question Apakah anda bisa melihat pertanyaan 2 ini?
\question Apakah anda bisa melihat pertanyaan 3 ini?
\question Apakah anda bisa melihat pertanyaan 4 ini?
\question Berikut adalah sub pertanyaan:
\begin{tasks}
\task Sub pertanyaan 1?
\task Sub pertanyaan 2?
\end{tasks}
\question Apakah anda bisa melihat pertanyaan 5 ini?
\end{questions}

\section{Logo Penerbit Jala Marwita Lestari}
\lipsum[3] Logo bisa dilihat pada gambar~\ref{fig:logo}.

\begin{figure}[h]
    \centering
    \includegraphics[scale=0.35]{img/logo-penerbit-jamali.png}
    \caption{Logo dari Penerbit Jala Marwita Lestari}
    \label{fig:logo}
\end{figure}

\chapter{Ini Bab Berikutnya}

Kali ini, kami akan menampilkan contoh tabel pada Tabel~\ref{table:contoh-tabel}

{\rowcolors{3}{green!80!yellow!50}{green!70!yellow!40}
\begin{table}[H]
\caption{Contoh Tabel}
\label{table:contoh-tabel}
\begin{tabular}{ |p{5cm}|p{3cm}|p{3cm}|  }
\hline
\multicolumn{3}{|c|}{Country List} \\
\hline
Country Name or Area Name& ISO ALPHA 2 Code &ISO ALPHA 3 \\
\hline
Afghanistan & AF &AFG \\
Aland Islands & AX   & ALA \\
Albania &AL & ALB \\
Algeria    &DZ & DZA \\
American Samoa & AS & ASM \\
Andorra & AD & AND   \\
Angola & AO & AGO \\
\hline
\end{tabular}
\end{table}}

\section{Sub Bab Nih}

Ini adalah contoh menggunakan verb: \verb|Penerbit Jala Marwita Lestari|. \lipsum[4]

\section{Ini Sub Bab Juga}

Kali ini, sub bab akan diisi contoh source code.

\begin{lstlisting}[language=Python, caption=Contoh Kode Sumber Python]
import numpy as np
    
def incmatrix(genl1,genl2):
    m = len(genl1)
    n = len(genl2)
    M = None #to become the incidence matrix
    VT = np.zeros((n*m,1), int)  #dummy variable
    
    #compute the bitwise xor matrix
    M1 = bitxormatrix(genl1)
    M2 = np.triu(bitxormatrix(genl2),1) 

    for i in range(m-1):
        for j in range(i+1, m):
            [r,c] = np.where(M2 == M1[i,j])
            for k in range(len(r)):
                VT[(i)*n + r[k]] = 1;
                VT[(i)*n + c[k]] = 1;
                VT[(j)*n + r[k]] = 1;
                VT[(j)*n + c[k]] = 1;
                
                if M is None:
                    M = np.copy(VT)
                else:
                    M = np.concatenate((M, VT), 1)
                
                VT = np.zeros((n*m,1), int)
    
    return M
\end{lstlisting}

\end{document}
